\documentclass{ximera}

\input{../../preamble.tex}

\title{Core Functions}

\begin{document}

\begin{abstract}
parents
\end{abstract}
\maketitle




In this course, we will begin studying a collection of functions called the \textbf{Elementary Functions}.

The \textbf{Elementary Functions} are created from an initial collection of ``building block'' functions, which we will call the \textbf{Core Functions}.

The Elementary Functions are created by combining the Core Functions together through addition, subtraction, multiplication, division, composition, and inverses.


The Elementary Functions are sums, differences, products, quotients, compositions, and inverses of Core Functions.

Calculus is largly about analyzing the Elementary Functions.






We will begin with seven Core Functions

\begin{itemize}
\item{Constant Functions: $f(x) = c$}
\item{The Identity Function: $f(x) = x$}
\item{Power Functions: $f(x) = x^n$}
\item{Root Functions: $f(x) = \sqrt[n]{x}$}
\item{The Exponential Function: $f(x) = e^x$}
\item{The Natural LogarithmiFunction: $f(x) = \ln(x)$}
\item{Sine and Cosine: $f(x) = sin(x)$ and $f(x) = cos(x)$}
\item{The Absolute Value Function: $f(x) = |x|$}
\end{itemize}


By forming sums, differences, products, quotients, compositions, and inverses of these Core Functions, we obtain the Elementary Functions.


\begin{itemize}
\item{Constant}
\item{Linear}
\item{Quadratic}
\item{Polynomial}
\item{Rational}
\item{Radical/Root}
\item{Exponential}
\item{Shifted Exponential}
\item{Logarithmic}
\item{Absolute Value}
\item{Trigonometric}
\item{Inverse Trigonometric}
\end{itemize}


Creating the library of Elementray Functions is a long journey. We'll need a lot of experiencing, which will take us right through Calculus.







\subsection*{The Core Functions}











\begin{onlineOnly}
\begin{center}
\textbf{\textcolor{green!50!black}{ooooo-=-=-=-ooOoo-=-=-=-ooooo}} \\

more examples can be found by following this link\\ \link[More Examples of Formulas]{https://ximera.osu.edu/csccmathematics/precalculus1/precalculus1/formulas/examples/exampleList}

\end{center}
\end{onlineOnly}





\end{document}
